\section{A worked protocol for one elementary step}
\label{sec:worked}

Section~\ref{sec:intro} listed the inputs a certificate needs. This section
supplies them for one system far enough to expose where the method is already
informative and where it is not. Nothing below is a molecular energy
calculation; the certified content is arithmetic on published growth data plus
two geometric inequalities, and the negative result is proved.

\subsection{System and evidence status}

The target is $\mathrm{CaF_2}$ growth from a calcium amidinate and hydrogen
fluoride on the fluorine-terminated $\mathrm{CaF_2}(111)$ cleavage plane.
Every choice below is either taken from a cited measurement or is explicitly
labelled as a modelling decision.

\begin{itemize}
\item \emph{Surface.} The $(111)$ plane is the natural cleavage plane and is
F-terminated with a $(1\times1)$ unreconstructed termination
\citep{Vogt2005,Puchin2001}. Its repeat unit along $[111]$ is a
F--Ca--F trilayer of height $a_0/\sqrt3$ \citep{Calleja2005}.
\item \emph{Lattice constant.} $a_0 = 5.4626$~\AA, room-temperature XRD powder
data (Swanson \& Tatge, NBS Circular 539 (1953) 69--70; ICSD 52754).
\item \emph{Precursor pair.} Bis($N,N'$-di-\emph{i}-propylformamidinato)calcium(II)
with anhydrous HF, reported at 175--250~$^\circ$C with a growth per cycle of
$0.3$--$0.4$~\AA/cycle \citep{Hennessy2023}.
\item \emph{Competing routes in the literature.} $\mathrm{Ca(thd)_2}$ with
$\mathrm{TiF_4}$ at 300--450~$^\circ$C \citep{Pilvi2007,Pilvi2008}; HF--pyridine
as a fluorine co-reactant \citep{Lee2016}; fluorinated $\beta$-diketonates as
single-source precursors \citep{LoPresti2022}; the first metal-fluoride ALD
process \citep{Ylilammi1994}.
\end{itemize}

The literature does not fix the elementary-step mechanism. Whether HF
protonates an adsorbed amidinate, exchanges ligands, or first builds a surface
fluorine reservoir is an open question, and no activation energy for this
system has been measured or computed in any source consulted here. We therefore
certify only what survives that ignorance.

\subsection{Certified arithmetic from published growth data}

Two published numbers, one measured and one geometric, determine the fraction
of a layer deposited per cycle. Let $p$ be the trilayer height and $s$ the
growth per cycle. Writing $s = p\,c$ with $c$ the fractional coverage, an
interval $s\in[s_1,s_2]$ with $p>0$ gives $c\in[s_1/p,\ s_2/p]$
(Lemma~\ref{lem:coverage-lean}).

\begin{lemma}[Interval coverage]
\label{lem:coverage-lean}
For real $p>0$ and $s=p\,c$, if $s_1\le s\le s_2$ then
$s_1/p\le c\le s_2/p$.
\end{lemma}
\begin{proof}
From $s=p\,c$ we get $c=s/p$. Since $p>0$, division by $p$ is monotone, so
$s_1/p\le s/p\le s_2/p$.
\end{proof}

With $p=a_0/\sqrt3=3.1538$~\AA{} and $s\in[0.3,0.4]$~\AA/cycle,
\[
c \in \left[\frac{0.3}{3.1538},\ \frac{0.4}{3.1538}\right]
  = [9.51\%,\ 12.68\%].
\]
Each cycle therefore completes about one tenth of a trilayer.

On the unreconstructed F-terminated $(111)$ plane the Ca ions form a triangular
lattice of spacing $a_0/\sqrt2$, so one Ca site occupies
\[
A_{\mathrm{site}} = \tfrac{\sqrt3}{2}\left(\tfrac{a_0}{\sqrt2}\right)^2
= \tfrac{\sqrt3}{4}a_0^2 = 12.9211\ \text{\AA}^2 .
\]
The areal density of deposited Ca follows as
$\sigma = c/A_{\mathrm{site}} \in [0.00736,\ 0.00982]$~\AA$^{-2}$, i.e.\ one
deposited Ca per $101.9$--$135.8$~\AA$^2$ of surface. That is $7.9$ to $10.5$
times the area of a bare Ca site.

\begin{remark}
This ratio is the quantitative statement that the cycle is \emph{not} limited
by the availability of surface sites. It does not identify the limiting
mechanism.
\end{remark}

\subsection{A geometric design rule for the ligand}

The ratio above invites a purely geometric bound, which holds regardless of
kinetics.

\begin{theorem}[Footprint bound]
\label{thm:footprint}
Let $\sigma>0$ be the number of precursor molecules deposited per unit area and
let $A>0$ be the area each occupies, with no overlap. Then $\sigma A\le1$, and
if $A_0\le A$ for some ligand footprint $A_0>0$ then
\[
\sigma \le \frac{1}{A_0}.
\]
Equivalently, with growth per cycle $s = p\,\sigma A_{\mathrm{site}}$,
\[
s \le \frac{p\,A_{\mathrm{site}}}{A_0}.
\]
\end{theorem}
\begin{proof}
Non-overlap gives $\sigma A\le1$. Dividing by $\sigma>0$ yields
$A\le1/\sigma$; since $A_0\le A$ and both are positive,
$1/A\le1/A_0$, hence $\sigma\le1/A\le1/A_0$. Multiplying
$s=p\,\sigma A_{\mathrm{site}}$ by this bound gives the second inequality.
\end{proof}

The theorem is a cap, not a prediction: it constrains what a cycle can deposit
and says nothing about whether the cycle reaches the cap. Read backwards,
however, it is a design rule. A ligand of footprint $A_0$ caps the growth per
cycle at $pA_{\mathrm{site}}/A_0$, so a bulky ligand cannot produce a high
growth rate no matter how favourable its surface chemistry. For the present
system the observed $\sigma$ requires only $A_0\le135.8$~\AA$^2$, which is
consistent with, but does not establish, steric self-limitation by the
di-\emph{i}-propylformamidinate ligand. A kinetic limitation would produce the
same observation.

\subsection{Why sampling cannot certify a barrier}

It is tempting to sample energies at a set of candidate configurations and read
a barrier off their maximum. That inference is not available, and the failure
is elementary.

\begin{theorem}[No upper bound from sampling]
\label{thm:nosampling}
There is a configuration set $X$ with a subset $S\subseteq X$, nonempty, and an
energy function $E:X\to\R$ such that $E=0$ on all of $S$ while
$\max_X E = 10$. Consequently no function of the sampled values
$\{E(x):x\in S\}$ alone bounds $\max_X E$ from above.
\end{theorem}
\begin{proof}
Take $X=\{0,1,2\}$, $S=\{0,2\}$, and $E(0)=E(2)=0$, $E(1)=10$. Then
$\max_S E = 0$ and $\max_X E = 10$, and the sampled data are identical for this
$E$ and for the function that is identically zero, which has maximum $0$. Any
claimed upper bound derived from the samples would therefore have to be at
least $10$ and at most $0$.
\end{proof}

The positive direction survives: a sampled lower bound is valid at every
configuration where it is evaluated, so its maximum over the samples is a valid
lower bound on any barrier whose path passes through those configurations
(Lemma~\ref{lem:sample-lower}).

\begin{lemma}[Lower bound from samples]
\label{lem:sample-lower}
Let $L(x)\le E(x)$ for $x\in S$ and let $E_r\le U_r$. Then for every $x\in S$,
$L(x)-U_r \le E(x)-E_r$.
\end{lemma}
\begin{proof}
Immediate from $L(x)\le E(x)$ and $-U_r\le-E_r$.
\end{proof}

The asymmetry is the reason a certificate must treat the two directions
separately. When the interval produced for a barrier is one-sided, it must be
reported as one-sided.

\subsection{Formal verification}

Lemmas~\ref{lem:coverage-lean} and \ref{lem:sample-lower} and
Theorems~\ref{thm:footprint} and \ref{thm:nosampling} are formalized in Lean~4
against \texttt{mathlib} as \texttt{gpc\_coverage\_interval},
\texttt{barrier\_lower\_from\_samples},
\texttt{max\_density\_from\_ligand\_size} with
\texttt{gpc\_upper\_from\_ligand\_footprint}, and
\texttt{sampling\_gives\_no\_upper\_bound} respectively, together with the
supporting interval lemmas \texttt{adsorption\_energy\_interval},
\texttt{two\_corrections\_widen}, \texttt{coverage\_gpc\_interval},
\texttt{packing\_bound}, \texttt{barrier\_upper\_at\_max},
\texttt{certified\_barrier\_interval}, \texttt{stays\_positive} and
\texttt{certified\_negative}. Each is checked with \texttt{\#print axioms} and
depends only on \texttt{propext}, \texttt{Classical.choice} and
\texttt{Quot.sound}; no \texttt{sorry}, \texttt{admit} or project-defined axiom
appears in the development. Kernel acceptance certifies the inequalities above
given their hypotheses. It certifies nothing about the hypotheses themselves,
which remain physical or geometric inputs.

\subsection{What remains open for this system}

\begin{enumerate}
\item The elementary-step mechanism and any activation energy for
HF reacting with an adsorbed calcium amidinate on $\mathrm{CaF_2}(111)$. No
value is available in the sources consulted, and Theorem~\ref{thm:nosampling}
shows that a barrier can be certified from below only.
\item The ligand footprint $A_0$ of di-\emph{i}-propylformamidinate, which
Theorem~\ref{thm:footprint} needs in order to become a quantitative cap.
\item Whether the observed $9.51$--$12.68\%$ coverage reflects steric
self-limitation or kinetics. Distinguishing them requires a coverage-dependent
measurement, not a bound.
\item The substrate relationship, which is deliberately absent here. If growth
continues on existing $\mathrm{CaF_2}$ the relevant surface is $\mathrm{CaF_2}$;
the role of an underlying Si(100) is epitaxial orientation, not surface
chemistry.
\end{enumerate}

The honest summary of this section is therefore narrow. Two published numbers
plus a lattice constant fix the fractional coverage of a cycle and show that
site availability is not the limiter. A kernel-checked inequality turns that
observation into a ligand-design cap, and a kernel-checked counterexample shows
that the barrier for the same system cannot be bounded above by any amount of
configuration sampling. No barrier value for this system is claimed.
